\documentclass[11pt,a4paper]{article}

% ---------------------------------------------------------------------------
% Packages
% ---------------------------------------------------------------------------
\usepackage[utf8]{inputenc}
\usepackage[T1]{fontenc}
\usepackage[margin=1in]{geometry}
\usepackage{booktabs}
\usepackage{longtable}
\usepackage{array}
\usepackage{enumitem}
\usepackage[table]{xcolor}
\usepackage{titlesec}
\usepackage{amsmath}
\usepackage{amssymb}
\usepackage{graphicx} % Required for inserting architecture images
\usepackage[hidelinks]{hyperref}

% ===== Colors =====
\definecolor{lightblue1}{RGB}{235,245,255}
\definecolor{lightblue2}{RGB}{248,252,255}


\titleformat{\section}{\large\bfseries\color{black!85}}{\thesection}{0.6em}{}
\titleformat{\subsection}{\normalsize\bfseries}{\thesubsection}{0.6em}{}
\setlist{itemsep=2pt,topsep=3pt}
\renewcommand{\arraystretch}{1.25}

\newcommand{\dsname}[1]{\textbf{#1}}

% ---------------------------------------------------------------------------
\title{\vspace{-1.5em}\textbf{Generative-AI--Based Automated Analysis and\\
Risk Scoring of Fraudulent APKs}\\[0.4em]
\large Datasets and Solution Approach}
\author{}
\date{}

\begin{document}
\maketitle
\vspace{-2em}

\begin{center}
{\Large \textbf{Team Chernobyl}}

\vspace{0.8em}

{\large
Bismay Samal \\
Mohit Ranjan Naik \\
Amiya Behera \\
Shreya Mishra
}
\end{center}

\vspace{2em}

\begin{abstract}
\noindent
Mobile-banking fraud increasingly relies on malicious Android application packages
(APKs) that impersonate legitimate banking and payment applications to steal one-time
passwords, harvest credentials, abuse accessibility services, and exfiltrate sensitive
data. Manual analysis of such samples is accurate but slow and does not scale with
submission volume. This document specifies the datasets we will use to train and
validate the machine-learning component of an automated APK risk-assessment platform,
and the layered solution approach we will follow. The approach fuses transparent
rule-based scoring, supervised machine learning, controlled dynamic behaviour analysis,
and a generative-AI reasoning layer to produce explainable risk scores and
analyst-ready investigation reports, with explicit attention to concept drift,
false-positive control, and rigorous, honest evaluation.
\end{abstract}

\vspace{1em}


% ===========================================================================
\section{Introduction}
% ===========================================================================
Android remains the dominant mobile platform and, consequently, the primary target of
mobile-banking fraud. Malicious APKs increasingly impersonate legitimate banking,
payment, and utility applications to steal one-time passwords (OTPs), harvest
credentials, abuse accessibility services, and exfiltrate sensitive data. Manual triage
of such samples is slow and does not scale with the volume of submissions a fraud or
Security Operations Centre (SOC) team receives.

This document describes the \emph{datasets we will use} to train and validate the
machine-learning component of our platform, and the \emph{solution approach} we will
follow to deliver automated, explainable risk assessment of suspicious APKs. Our system
combines transparent rule-based scoring, supervised machine learning, controlled dynamic
behaviour analysis, and a generative-AI reasoning layer that produces human-readable
investigation reports.

\begin{figure}[htbp]
\centering

\begin{minipage}{0.22\textwidth}
\centering 
\includegraphics[width=\textwidth]{manual_apk_analysis_challenges.png}
\end{minipage}
\hspace{0.01\textwidth}
\begin{minipage}{0.22\textwidth}
\centering
\includegraphics[width=\textwidth]{mobile_fraud_threat_landscape.png}
\end{minipage}
\hspace{0.01\textwidth}
\begin{minipage}{0.22\textwidth}
\centering
\includegraphics[width=\textwidth]{apkguard_solution_overview.png}
\end{minipage}

\caption{Introduction overview illustrating the mobile fraud threat landscape, challenges of manual APK investigation, and the motivation for an automated AI-powered APK risk assessment platform.}
\label{fig:introduction_overview}
\end{figure}

\subsection{Objectives}
Our objectives are to: (i)~extract security-relevant evidence from an APK through static
and (optionally) dynamic analysis; (ii)~produce a calibrated, explainable risk score that
discriminates malicious from benign applications while keeping false positives low;
(iii)~train a machine-learning classifier that generalises to recent, previously unseen
samples; and (iv)~generate human-readable investigation reports that communicate findings
and recommended actions to both technical and non-technical stakeholders.

\subsection{Design principles}
Three principles guide the design. \emph{Explainability:} every risk contribution must be
traceable to concrete evidence. \emph{Safety:} potentially malicious code is executed only
inside an isolated, disposable environment, never on analyst infrastructure.
\emph{Honesty:} reported confidence reflects the evidence, and known limitations (drift,
evasion, unavailable C2) are stated rather than hidden.

% ===========================================================================
\section{Background}
% ===========================================================================
\subsection{Static versus dynamic analysis}
\emph{Static analysis} inspects an application without running it---parsing the manifest,
decompiling bytecode, and extracting permissions, API references, strings, and structural
signals. It is fast and safe but can be defeated by encryption, packing, and runtime code
loading. \emph{Dynamic analysis} executes the application in a controlled environment and
observes actual behaviour---network traffic, file activity, and security-sensitive API
calls---revealing intent that static analysis cannot, at the cost of greater complexity
and the need for isolation. The two are complementary, and our approach uses both.

\subsection{Machine learning for malware detection}
Supervised learning over engineered features (permissions, API usage, intents, network
indicators) is a well-established approach to Android-malware detection. Its principal
challenges are \emph{class imbalance} (benign samples dominate), \emph{concept drift}
(malware evolves, degrading models trained on older data), and \emph{evaluation inflation}
(near-duplicate samples leaking between train and test). Our methodology explicitly
addresses all three.

\subsection{Generative AI as a reasoning layer}
Large language models excel at correlating heterogeneous evidence and producing fluent,
structured explanations. In our design, generative AI does \emph{not} act as a detector;
it reasons over already-extracted findings to produce summaries, attack-chain narratives,
and recommendations, and to contextualise ambiguous evidence---for example, distinguishing
benign third-party library usage from genuinely malicious behaviour.

% ===========================================================================
\section{Datasets}
% ===========================================================================
The machine-learning model will be trained, validated, and stress-tested on a
combination of established Android-malware corpora and curated benign application sets.
We deliberately combine \emph{older, well-labelled} datasets (for breadth of malware
families and reproducibility) with \emph{recent} datasets and live sampling (to mitigate
concept drift, since malware behaviour evolves over time).

\subsection{Primary training and benchmark datasets}

\begin{itemize}
    \item \dsname{Drebin.} A widely cited Android-malware benchmark containing approximately \textbf{5{,}560 malicious} samples spanning \textbf{179 malware families}, alongside a large set of benign applications. Crucially, Drebin provides \emph{pre-extracted static feature sets} per sample---requested and used permissions, API calls, intent actions, hardware components, application components, and embedded network indicators---together with a family-label mapping. These ready-made feature vectors allow us to train and benchmark classifiers without reprocessing raw APKs, making it our primary baseline dataset.
    \item \dsname{CICMalDroid 2020.} A modern dataset of roughly \textbf{17{,}000} Android samples organised into five classes: \emph{Adware}, \emph{Banking malware}, \emph{SMS malware}, \emph{Riskware}, and \emph{Benign}. It includes both static attributes and dynamically observed behaviour (system calls, binder activity, composite behaviours). Its banking- and SMS-malware categories align directly with our fraud-triage objective, and its recency helps the model generalise beyond the older Drebin distribution.
    \item \dsname{CIC-AndMal2017.} A dataset emphasising \emph{dynamic and network-traffic} features, comprising several hundred malware samples across major families (ransomware, scareware, SMS malware, adware) and several thousand benign applications. We will use it to enrich the feature space with runtime and network-flow characteristics that static analysis cannot capture.
    \item \dsname{AndroZoo.} A continuously growing repository of \textbf{millions} of real-world APKs collected from Google Play and numerous third-party markets, each annotated with multi-engine anti-virus detection counts and collection metadata. We will use AndroZoo as a source of \emph{fresh, diverse} benign and malicious samples drawn from recent time windows, specifically to build an out-of-distribution evaluation set for measuring concept drift and real-world false-positive rates.
    \item \dsname{Android Malware Genome Project (MalGenome).} A historical reference set of \textbf{1{,}260} malware samples across \textbf{49 families}. Although dated, it provides a clean, well-documented baseline of classic malware behaviours that is useful for sanity-checking detectors and for family-level analysis.
    \item \dsname{KronoDroid.} A hybrid dataset that records \emph{timestamped} samples with both static and dynamic features captured on real devices and emulators. Its temporal labelling makes it particularly valuable for time-aware evaluation, i.e.\ training on earlier samples and testing on later ones to simulate realistic deployment and quantify drift.
\end{itemize}

\subsection{Benign application sources}
A robust classifier requires a representative, modern benign distribution to keep
false-positive rates low. In addition to the benign partitions of the datasets above, we
will curate benign applications from:
\begin{itemize}
    \item \dsname{F-Droid} --- a catalogue of free and open-source Android applications, useful because the sources are auditable and the apps are demonstrably non-malicious.
    \item \dsname{Google Play (top and long-tail apps via AndroZoo)} --- to capture the permission and API profiles of mainstream commercial applications, including advertisement- and analytics-heavy apps that are benign but superficially ``suspicious.''
\end{itemize}

\subsection{Threat-intelligence corroboration source}
For ground-truth verification and IOC enrichment (not as a training feature), we will
consult \dsname{VirusTotal}, which aggregates verdicts from many anti-malware engines.
This supports label validation and provides an independent second opinion during
evaluation.

\subsection{Feature taxonomy}
Features are grouped into static and dynamic families, as summarised in
Table~\ref{tab:features}. We prioritise features expressed \emph{consistently} across data
sources, so that a model trained offline applies directly to live submissions with minimal
train/serve skew.

\begin{longtable}{p{3.0cm} p{2.4cm} p{8.0cm}}
\toprule
\textbf{Group} & \textbf{Origin} & \textbf{Example features}\\
\midrule
\endhead
Permissions & Static & Requested/used permission vectors; counts of high-risk permissions (SMS, contacts, accessibility, overlay, install-packages)\\
API usage & Static & Presence/frequency of security-relevant API references (SMS sending, accessibility, dynamic code loading, command execution, networking)\\
Components \& intents & Static & Activities/services/receivers/providers; exported components; intent-filter actions\\
Network indicators & Static & Counts of URLs, IPs, domains; embedded tokens; decoded (Base64/hex) indicators; native-library strings\\
Obfuscation/packing & Static & Identifier-renaming ratio; multiple DEX files; high-entropy blobs; packer/compiler fingerprints\\
Certificate & Static & Debug/placeholder identity flags; validity anomalies\\
Runtime behaviour & Dynamic & System/API calls; started services; granted runtime permissions; file activity\\
Network flow & Dynamic & Contacted endpoints; protocol/port; traffic-flow statistics\\
\bottomrule
\caption{Feature taxonomy used for classification.}
\label{tab:features}
\end{longtable}

\subsection{Summary of Datasets}

\rowcolors{2}{lightblue1}{lightblue2}

\begin{longtable}{|p{2.8cm}|p{2.0cm}|p{3.0cm}|p{5.5cm}|}
\caption{Datasets used for model training, benchmarking, and evaluation.} \\
\hline
\rowcolor{blue!25}
\textbf{Dataset} &
\textbf{Approx. Scale} &
\textbf{Feature Type} &
\textbf{Role in Our Work} \\
\hline
\endfirsthead

\hline
\rowcolor{blue!25}
\textbf{Dataset} &
\textbf{Approx. Scale} &
\textbf{Feature Type} &
\textbf{Role in Our Work} \\
\hline
\endhead

Drebin &
5.5k malware + benign &
Static (permissions, APIs, intents, IOCs) &
Primary baseline dataset with pre-extracted features for rapid model training.
\\ \hline

CICMalDroid 2020 &
$\sim$17k &
Static + Dynamic &
Modern multi-class dataset including banking and SMS malware.
\\ \hline

CIC-AndMal2017 &
$\sim$5.5k &
Dynamic + Network &
Provides runtime and network-flow feature enrichment.
\\ \hline

AndroZoo &
Millions &
Raw APKs + AV Labels &
Fresh real-world APK sampling for concept drift and false-positive evaluation.
\\ \hline

MalGenome &
1.26k &
Static &
Historical malware benchmark for family-level validation.
\\ \hline

KronoDroid &
Tens of Thousands &
Hybrid, Timestamped &
Supports temporal and concept-drift evaluation.
\\ \hline

F-Droid / Google Play &
Curated &
Raw APKs &
Modern benign application distribution.
\\ \hline

VirusTotal &
On-demand &
Multi-engine Verdicts &
Label validation, IOC corroboration, and reputation verification.
\\ \hline

\end{longtable}


\subsection{Data preparation and partitioning}
\begin{itemize}
    \item \textbf{Normalisation and labelling.} Sample identities will be tracked by cryptographic hash; labels will be reconciled across datasets and corroborated with multi-engine verdicts to reduce label noise.
    \item \textbf{Class balance.} Because benign samples vastly outnumber malicious ones, we will apply class weighting and/or balanced resampling to avoid a model biased toward the majority class.
    \item \textbf{Splitting strategy.} We will use stratified train/validation/test splits for in-distribution evaluation and, additionally, a \emph{temporal split} (train on older samples, test on newer ones) to obtain an honest estimate of real-world performance under concept drift.
    \item \textbf{Leakage control.} Family-aware and time-aware splitting will be used to prevent near-duplicate samples from the same campaign appearing in both training and test sets, which would otherwise inflate reported accuracy.
\end{itemize}

% ===========================================================================
\section{Solution Approach}
% ===========================================================================
Our approach treats APK risk assessment as a layered pipeline. Each layer contributes
independent, explainable evidence, and the layers are fused into a single, transparent
risk score accompanied by a natural-language report. Generative AI is used to
\emph{reason over and communicate} the extracted evidence, never as the sole detector.

\begin{figure}[htbp]
    \centering
    % Make sure "WhatsApp Image 2026-06-15 at 12.43.27 PM.jpeg" matches the image filename in your project folder
    \includegraphics[width=0.85\textwidth]{WhatsApp Image 2026-06-15 at 12.43.27 PM.jpeg}
    \caption{Layered processing pipeline spanning sandbox instrumentation, dynamic behavior tracking, feature extraction, and ingestion into the machine learning input layer.}
    \label{fig:pipeline_architecture}
\end{figure}

\subsection{Stage 1 --- Intake and identity}
Each submitted APK will be hashed (e.g.\ SHA-256/MD5) for de-duplication and blacklisting,
and basic metadata (package name, version, target/minimum SDK, file size) will be
recorded. Hash-based identity enables consistent labelling and correlation with
threat-intelligence feeds.

\subsection{Stage 2 --- Static analysis}
Static analysis inspects the application without executing it and forms the backbone of
fast, safe triage. We will extract:
\begin{itemize}
    \item \textbf{Manifest and permissions:} declared components, intent filters, and the set of requested permissions, with emphasis on fraud-relevant permissions (SMS, contacts, accessibility, overlay, install-packages, phone state).
    \item \textbf{Suspicious API usage:} references to APIs associated with OTP theft, accessibility abuse, dynamic code loading, command execution, and network communication, weighted by how discriminative each is in benign versus malicious apps.
    \item \textbf{Indicators of compromise (IOCs):} URLs, IP addresses, domains, embedded tokens, and likely secrets, including indicators recovered by decoding obfuscated (e.g.\ Base64/hex) strings and by scanning bundled native libraries.
    \item \textbf{Obfuscation and packing signals:} identifier-renaming ratios, multiple DEX files, high-entropy blobs, dynamic loaders, and packer/compiler fingerprints.
    \item \textbf{Signing certificate and attack surface:} certificate red-flags (debug/placeholder identities, implausible validity) and exported components reachable by other applications.
\end{itemize}
Signature-based screening (rule sets expressing known fraud patterns) and a mapping of
findings to a standard adversary-technique taxonomy (MITRE ATT\&CK for Mobile) will be
applied to make the output actionable for analysts.

A central design concern is \emph{false-positive control}. Many indicators that appear
alarming in isolation---requesting network access, referencing cryptography or reflection
APIs, or bundling advertising and analytics libraries---are common in entirely legitimate
applications. Accordingly, each static signal is weighted by how \emph{discriminative} it
is between benign and malicious applications: rare, high-signal capabilities (programmatic
SMS sending, accessibility automation, runtime code loading) carry substantial weight,
whereas ubiquitous, low-signal API references are weighted near zero. Standard release-
engineering practices (identifier minification, common packers and compilers) are
recognised and not penalised as if they were evasion. This calibration is what allows the
system to flag genuine fraud while keeping mainstream, advertisement-heavy applications
out of the high-severity bands.

\subsection{Stage 3 --- Controlled dynamic analysis}
Because sophisticated samples conceal behaviour from static inspection (e.g.\ encrypted
payloads decrypted only at runtime), we will optionally execute samples inside an
\emph{isolated, instrumented sandbox}. Many fraudulent applications ship no launcher
icon and instead activate through broadcast receivers (such as device-boot or
connectivity events) and background services; a naive ``launch the main activity''
strategy would therefore observe nothing. To address this, the sandbox systematically
\emph{exercises every declared entry point}---the launcher activity, each broadcast
receiver (driven with its registered intent actions), and each declared service---and
automatically dismisses first-run permission and consent dialogs so that execution is not
stalled behind a modal prompt. It then records runtime behaviour: granted permissions,
started services, contacted network endpoints, file-system activity, and
security-relevant API invocations.

To counter anti-analysis behaviour, the sandbox applies \emph{evasion-resistance}. Many
samples fingerprint the environment (device build properties, system properties,
telephony identifiers, and the presence of emulator-specific files) and remain dormant
when they detect analysis. The sandbox presents realistic device characteristics so that
these emulator-detection checks pass and the genuine malicious logic executes. Through
dynamic instrumentation it further intercepts security-relevant calls at runtime to
\emph{capture decrypted second-stage payloads and command-and-control (C2) endpoints}
that never appear in the static image---high-value indicators that would otherwise be
invisible. All execution is confined to a disposable, network-restricted, isolated
environment with snapshot restoration, and is strictly opt-in; the static pipeline never
executes the sample.

\subsection{Stage 4 --- Feature engineering}
Static and dynamic observations will be transformed into a fixed-length feature
representation for the classifier. The representation will be drawn from features that
are expressed consistently across data sources, including:
\begin{itemize}
    \item permission presence vectors;
    \item suspicious-API presence/frequency vectors;
    \item intent and component indicators;
    \item counts of network indicators and obfuscation signals;
    \item selected runtime-behaviour and network-flow features (when dynamic analysis is available).
\end{itemize}
We will favour features that are robust to superficial obfuscation and that remain
comparable between training data and live submissions, in order to minimise feature
drift between offline training and deployment.

\begin{figure}[htbp]
\centering
\includegraphics[width=0.92\textwidth]{solution_approach.png}
\caption{Solution approach workflow of APKGuard AI. The pipeline begins with APK intake and identity management, followed by static and dynamic analysis, feature engineering, machine-learning classification, hybrid risk scoring, threat-intelligence corroboration, and evidence-grounded Generative AI reporting to produce actionable security assessments.}
\label{fig:solution_approach}
\end{figure}


\subsection{Stage 5 --- Machine-learning classification}
We will train supervised classifiers to estimate the probability that a sample is
malicious and, where data permits, its likely family/category (e.g.\ banking, SMS,
adware). Candidate models include:
\begin{itemize}
    \item \textbf{Gradient-boosted trees} (e.g.\ XGBoost/LightGBM) and \textbf{Random Forests}, chosen for strong tabular-data performance and interpretability through feature-importance analysis;
    \item \textbf{Logistic Regression} as a transparent baseline;
    \item Optional advanced directions, such as graph-based models over the application call graph, as a future extension.
\end{itemize}
The classifier output is treated as one \emph{corroborating signal} fused with the
rule-based score---not as a stand-alone verdict---so that the final assessment remains
explainable.

Training follows a disciplined protocol. Model selection and hyperparameter tuning use
cross-validation on the training partition only; the held-out and temporal test sets are
touched once, at the end, to obtain an unbiased estimate. Class imbalance is handled with
class weighting and/or balanced resampling, and predicted probabilities are
\emph{calibrated} (e.g.\ via isotonic or Platt scaling) so that a reported probability of,
say, $0.8$ corresponds to a genuine $80\%$ likelihood---a prerequisite for fusing the ML
score meaningfully with the rule-based score. Where category labels are available
(banking, SMS, adware, riskware), we additionally train a multi-class head to suggest the
likely malware type, which enriches the analyst report. Feature importance and per-sample
attribution are retained so that each ML-driven decision can be explained in terms of the
specific permissions and APIs that drove it.

\subsection{Stage 6 --- Hybrid risk scoring}
The final risk score will combine the transparent rule-based score, the machine-learning
probability, and threat-intelligence corroboration:
\[
\mathrm{Risk} \;=\; \mathrm{RuleScore} \;+\; \alpha\,\mathrm{ML} \;+\; \beta\,\mathrm{Intel},
\]
normalised to a $0$--$100$ scale and mapped to the severity bands of
Table~\ref{tab:severity}. Every contribution is itemised so an analyst can see exactly
\emph{why} a sample received its score. Confirmed multi-engine detections raise
confidence, while patterns characteristic of benign advertisement-heavy applications are
down-weighted to control false positives. The fusion weights $\alpha,\beta$ are calibrated
on the validation set to optimise the precision/recall trade-off at the chosen operating
point, so that neither the ML signal nor external intelligence can override the
explainable rule-based evidence on its own.

\begin{table}[h]
\centering
\begin{tabular}{ll}
\toprule
\textbf{Score range} & \textbf{Severity}\\
\midrule
0--20 & Low\\
21--40 & Medium\\
41--70 & High\\
71--100 & Critical\\
\bottomrule
\end{tabular}
\caption{Risk-score to severity mapping.}
\label{tab:severity}
\end{table}

\subsection{Stage 7 --- Generative-AI reasoning and reporting}
A generative-AI layer will consume only the \emph{extracted findings} (never the raw
application) and produce: an executive summary for non-technical fraud managers, a
technical narrative correlating the evidence into a likely attack chain, a list of
indicators of compromise, and recommended actions for fraud/SOC teams. The model is
prompted to reason critically about the evidence---distinguishing, for instance, benign
library usage from genuine malicious intent---so that the report adds analytical value
rather than merely restating the findings.

To preserve trust, the layer is \emph{evidence-grounded}: the model is constrained to the
structured findings supplied to it and is instructed not to introduce indicators or
claims that are not present in that evidence, which suppresses hallucination. The output
is requested in a fixed, structured schema so that it can be rendered consistently in the
dashboard and exported into the formal report. Sending only distilled findings---rather
than the raw application---also keeps the design safe and privacy-preserving, since no
potentially malicious artefact is transmitted to the language model. Should the
generative service be unavailable, a deterministic template produces an equivalent report
from the same findings, so the platform degrades gracefully.

As a controlled extension, this reasoning layer will support a \emph{bounded agentic
investigation loop}. A local Ollama model (initially Qwen3 8B) may select and sequence
approved read-only operations such as manifest inspection, extracted-string search, and
bounded call-graph queries. Tool calls are schema-validated, evidence-cited, subject to
step and time budgets, and recorded in an audit trace. Gemini remains an optional
reviewer and report writer over the same redacted evidence bundle. External-intelligence
queries and emulator-based dynamic analysis require explicit operator approval; the
agent is never given a shell, arbitrary network access, or permission to execute an APK
on the host. Neither model assigns or changes the risk score, which remains deterministic
and explainable.

\subsection{Stage 8 --- Output and integration}
The platform will present an analyst dashboard (risk gauge, itemised breakdown, evidence
panels, and the AI report) and will export machine-readable artefacts: a formal PDF
investigation report and IOC lists suitable for ingestion into blocklists and SOC
tooling.
\subsection{Platform Architecture}

The APKGuard AI platform follows a layered architecture that integrates APK ingestion, static and dynamic analysis engines, feature engineering pipelines, machine-learning classifiers, threat-intelligence services, hybrid risk-scoring mechanisms, and Generative AI reporting components. Each layer contributes explainable evidence that is progressively refined into actionable security intelligence for analysts and fraud investigation teams.

\begin{figure}[htbp]
\centering
\includegraphics[width=\textwidth]{platform_architecture.png}
\caption{High-level platform architecture of APKGuard AI showing the end-to-end flow from APK submission and analysis through feature engineering, machine-learning classification, hybrid risk scoring, and Generative AI powered reporting.}
\label{fig:platform_architecture}
\end{figure}

\subsection{Technology Stack Map}

The implementation of APKGuard AI relies on a modern technology ecosystem spanning frontend development, backend orchestration, malware-analysis frameworks, machine-learning libraries, Generative AI services, databases, cloud infrastructure, monitoring systems, and deployment tooling. This technology stack enables scalable, secure, and maintainable operation of the platform.

\begin{figure}[htbp]
\centering
\includegraphics[width=\textwidth]{tech_stack_map.png}
\caption{Technology Stack Map illustrating the software frameworks, analysis tools, machine-learning libraries, databases, cloud services, DevOps infrastructure, and monitoring components used to implement the APKGuard AI platform.}
\label{fig:tech_stack_map}
\end{figure}



% ===========================================================================
\section{Evaluation Methodology}
% ===========================================================================
\subsection{Metrics}
Let TP, FP, TN, FN denote true/false positives/negatives. We report
\[
\mathrm{Precision}=\frac{\mathrm{TP}}{\mathrm{TP}+\mathrm{FP}},\quad
\mathrm{Recall}=\frac{\mathrm{TP}}{\mathrm{TP}+\mathrm{FN}},\quad
\mathrm{F1}=\frac{2\,\mathrm{Precision}\cdot\mathrm{Recall}}
{\mathrm{Precision}+\mathrm{Recall}},
\]
together with ROC-AUC and the false-positive rate
$\mathrm{FPR}=\mathrm{FP}/(\mathrm{FP}+\mathrm{TN})$. We emphasise precision/recall and
FPR over raw accuracy, since class imbalance makes accuracy misleading.

\subsection{Protocols}
\begin{itemize}
    \item \textbf{In-distribution:} stratified hold-out and cross-validation on each corpus.
    \item \textbf{Temporal (drift):} train on earlier samples, test on later ones, to estimate deployment performance as malware evolves.
    \item \textbf{Cross-dataset:} train on one corpus, evaluate on another, to measure generalisation.
    \item \textbf{False-positive focus:} dedicated evaluation against a curated set of modern benign applications, since over-flagging legitimate apps erodes analyst trust.
    \item \textbf{Explainability checks:} verify that high-scoring decisions are supported by human-interpretable evidence and that feature importances are consistent with domain knowledge.
\end{itemize}

\subsection{Baselines and ablations}
We compare the hybrid system against (i)~rule-based scoring alone and (ii)~the ML model
alone, and we ablate feature groups (permissions only, $+$API, $+$dynamic) to quantify the
marginal value of each evidence source and of dynamic analysis.

% ===========================================================================
\section{Threats to Validity and Limitations}
% ===========================================================================
\begin{itemize}
    \item \textbf{Concept drift.} Models trained on historical data degrade over time; temporal evaluation and periodic retraining mitigate but do not eliminate this.
    \item \textbf{Evasion.} Heavily packed samples, native-code logic, and anti-analysis checks can resist both static and dynamic inspection; evasion-resistance reduces but cannot fully remove this risk.
    \item \textbf{Unavailable infrastructure.} Samples whose payloads are fetched from now-defunct C2 servers may not exhibit full behaviour dynamically; in such cases the static evidence and captured indicators remain the authoritative signal.
    \item \textbf{Label noise.} Multi-engine verdicts disagree; threshold policies and reconciliation mitigate this, but residual uncertainty remains.
    \item \textbf{Dataset bias.} Public corpora over-represent certain families and eras; combining datasets with fresh sampling reduces, but does not remove, bias.
\end{itemize}

% ===========================================================================
\section{Ethical and Safety Considerations}
% ===========================================================================
All potentially malicious samples will be handled and executed only within isolated
environments using no personal accounts or networks. The system is designed for
\emph{defensive} purposes---first-level triage and investigation support for banking
fraud and SOC teams. Reporting will be faithful to the evidence, and limitations
(e.g.\ concept drift, sandbox evasion, payloads gated behind unavailable infrastructure)
will be stated transparently rather than overstated.

\begin{figure}[htbp]
\centering

% =========================
% GitHub Header
% =========================

\includegraphics[width=0.08\textwidth]{github_logo.png}

\vspace{0.5em}

{\LARGE\bfseries Source Code Repository}

\vspace{0.5em}

\href{https://github.com/IntrosWorld/malware-apk-detect}
{\Large\texttt{github.com/IntrosWorld/malware-apk-detect}}

\vspace{1.5em}

% =========================
% Two Column Layout
% =========================

\begin{minipage}[c]{0.42\textwidth}

\raggedright

{\Large\bfseries APKGuard AI Prototype}

\vspace{1em}

The APKGuard AI prototype demonstrates the complete end-to-end workflow of the proposed platform, integrating APK ingestion, static and dynamic malware analysis, feature engineering, machine-learning driven risk assessment, hybrid risk scoring, threat-intelligence enrichment, and Generative AI powered investigation reporting.

\vspace{0.8em}

The platform enables analysts to upload suspicious Android applications, inspect extracted indicators of compromise, visualize explainable risk scores, and generate evidence-grounded security reports for fraud investigation and SOC operations.

\vspace{1em}

{\large\bfseries Open Source • Research Prototype • Production Ready}

\end{minipage}
\hfill
\begin{minipage}[c]{0.53\textwidth}

\centering
\includegraphics[width=\textwidth]{demo.png}

\end{minipage}

\caption{APKGuard AI implementation and prototype demonstration. The GitHub repository provides the complete source code, while the prototype showcases automated APK analysis, threat detection, explainable risk scoring, and Generative AI powered reporting capabilities.}

\label{fig:prototype_demo}

\end{figure}

\newpage



% ===========================================================================
\section{Initial Prototype Photos}
% ===========================================================================

\begin{figure}[htbp]
    \centering
    \includegraphics[
        width=\linewidth,
        height=0.75\textheight,
        keepaspectratio
    ]{photo1.jpg}
    \caption{Prototype View 1}
    \label{fig:prototype_view_1}
\end{figure}

\begin{figure}[htbp]
    \centering
    \includegraphics[
        width=\linewidth,
        height=0.75\textheight,
        keepaspectratio
    ]{photo2.jpg}
    \caption{Prototype View 2}
    \label{fig:prototype_view_2}
\end{figure}

\begin{figure}[htbp]
    \centering
    \includegraphics[
        width=\linewidth,
        height=0.75\textheight,
        keepaspectratio
    ]{photo3.jpg}
    \caption{Prototype View 3}
    \label{fig:prototype_view_3}
\end{figure}

\begin{figure}[htbp]
    \centering
    \includegraphics[
        width=\linewidth,
        height=0.75\textheight,
        keepaspectratio
    ]{photo4.jpg}
    \caption{Prototype View 4}
    \label{fig:prototype_view_4}
\end{figure}

\begin{figure}[htbp]
    \centering
    \includegraphics[
        width=\linewidth,
        height=0.75\textheight,
        keepaspectratio
    ]{photo5.jpg}
    \caption{Prototype View 5}
    \label{fig:prototype_view_5}
\end{figure}


\end{document}
